% Options for packages loaded elsewhere
\PassOptionsToPackage{unicode}{hyperref}
\PassOptionsToPackage{hyphens}{url}
\PassOptionsToPackage{dvipsnames,svgnames,x11names}{xcolor}
%
\documentclass[
  10pt,
  a4paper,
]{article}
\usepackage{amsmath,amssymb}
\usepackage{setspace}
\usepackage{iftex}
\ifPDFTeX
  \usepackage[T1]{fontenc}
  \usepackage[utf8]{inputenc}
  \usepackage{textcomp} % provide euro and other symbols
\else % if luatex or xetex
  \usepackage{unicode-math} % this also loads fontspec
  \defaultfontfeatures{Scale=MatchLowercase}
  \defaultfontfeatures[\rmfamily]{Ligatures=TeX,Scale=1}
\fi
\usepackage{lmodern}
\ifPDFTeX\else
  % xetex/luatex font selection
  \setmainfont[]{Liberation Serif}
  \setsansfont[]{Liberation Sans}
  \setmonofont[]{DejaVu Sans Mono}
\fi
% Use upquote if available, for straight quotes in verbatim environments
\IfFileExists{upquote.sty}{\usepackage{upquote}}{}
\IfFileExists{microtype.sty}{% use microtype if available
  \usepackage[]{microtype}
  \UseMicrotypeSet[protrusion]{basicmath} % disable protrusion for tt fonts
}{}
\makeatletter
\@ifundefined{KOMAClassName}{% if non-KOMA class
  \IfFileExists{parskip.sty}{%
    \usepackage{parskip}
  }{% else
    \setlength{\parindent}{0pt}
    \setlength{\parskip}{6pt plus 2pt minus 1pt}}
}{% if KOMA class
  \KOMAoptions{parskip=half}}
\makeatother
\usepackage{xcolor}
\usepackage[margin=22mm]{geometry}
\setlength{\emergencystretch}{3em} % prevent overfull lines
\providecommand{\tightlist}{%
  \setlength{\itemsep}{0pt}\setlength{\parskip}{0pt}}
\setcounter{secnumdepth}{-\maxdimen} % remove section numbering
\ifLuaTeX
\usepackage[bidi=basic]{babel}
\else
\usepackage[bidi=default]{babel}
\fi
\babelprovide[main,import]{korean}
\ifPDFTeX
\else
\babelfont{rm}[]{Liberation Serif}
\fi
% get rid of language-specific shorthands (see #6817):
\let\LanguageShortHands\languageshorthands
\def\languageshorthands#1{}
\usepackage{xeCJK}
\setCJKmainfont{Noto Serif CJK KR}
\setCJKsansfont{Noto Sans CJK KR}
\setCJKmonofont{Noto Sans Mono CJK KR}
\xeCJKsetup{CJKspace=true}
\usepackage{xurl}
\usepackage{fvextra}
\usepackage{fancyhdr}
\AtBeginDocument{\DefineVerbatimEnvironment{verbatim}{Verbatim}{breaklines=true,breakanywhere=true,fontsize=\small}}
\setlength{\emergencystretch}{4em}
\setlength{\headheight}{15pt}
\pagestyle{fancy}\fancyhf{}
\fancyhead[L]{\small 중성자별 구조와 비방사 진동}
\fancyfoot[C]{\thepage}
\renewcommand{\headrulewidth}{0.3pt}
\setcounter{tocdepth}{2}
\ifLuaTeX
  \usepackage{selnolig}  % disable illegal ligatures
\fi
\usepackage{bookmark}
\IfFileExists{xurl.sty}{\usepackage{xurl}}{} % add URL line breaks if available
\urlstyle{same}
\hypersetup{
  pdftitle={A14 · Material interface의 전달과 restart},
  pdflang={ko-KR},
  colorlinks=true,
  linkcolor={Maroon},
  filecolor={Maroon},
  citecolor={Blue},
  urlcolor={Blue},
  pdfcreator={LaTeX via pandoc}}

\title{A14 · Material interface의 전달과 restart}
\usepackage{etoolbox}
\makeatletter
\providecommand{\subtitle}[1]{% add subtitle to \maketitle
  \apptocmd{\@title}{\par {\large #1 \par}}{}{}
}
\makeatother
\subtitle{과제 큐시트}
\author{}
\date{}

\begin{document}
\maketitle

\setstretch{1.12}
\textbf{선행:} A13, A09　 \textbf{독립 작업:} 4--7시간 (계획치)

\subsection{밀도가 뛰면 변위의 두 성분도 함께 뛰어야
하는가?}\label{uxbc00uxb3c4uxac00-uxb6f0uxba74-uxbcc0uxc704uxc758-uxb450-uxc131uxbd84uxb3c4-uxd568uxaed8-uxb6f0uxc5b4uxc57c-uxd558uxb294uxac00}

계면의 압력과 반경은 같고 에너지밀도만 다르다. 연속으로 둘 양을 평형
곡선이 아니라 채택한 유체 경계조건에서 먼저 고르자.

먼저 한두 문장으로 답한 뒤 아래 설명을 읽자. 첫 예측이 틀렸거나 아직
모르겠다고 적는 것은 감점하지 않는다.

\subsection{필요한 관계와 읽을
위치}\label{uxd544uxc694uxd55c-uxad00uxacc4uxc640-uxc77duxc744-uxc704uxce58}

\[
q_h=\frac{\epsilon_{\rm in}+P_t}{\epsilon_{\rm out}+P_t},
\qquad a=(\omega r_t/c)^2e^{\lambda_t-2\nu_t},
\] \[
W_{\rm out}=W_{\rm in},\qquad
V_{\rm out}=q_hV_{\rm in}+(q_h-1)\frac{\nu'_tW_{\rm in}}a.
\] \([\Delta P]=0\)과 zero-jump 한계를 직접 확인한다.

\begin{itemize}
\tightlist
\item
  \path|notes/equations/interface.md|
\item
  \path|learner/src/modes/ns_interface_transfer.f90|
\item
  \path|notes/chapters/ch07.md|
\end{itemize}

\subsection{계산으로
옮기기}\label{uxacc4uxc0b0uxc73cuxb85c-uxc62euxae30uxae30}

기본 경로는 제공 함수의 호출과 해석이다. 아래 \texttt{learner} 시험은
같은 물리 함수를 직접 작성한 경우에 실행한다. 제공 예제를 검사할 때에는
명시적으로 \texttt{calculation}을 선택한다.

\begin{enumerate}
\def\labelenumi{\arabic{enumi}.}
\tightlist
\item
  ΔP 식에서 enthalpy ratio를 직접 유도한다.
\item
  같은 r\_t의 inner/outer record를 따로 보존한다.
\item
  event에서 정지→transfer→restart를 수행한다.
\item
  zero jump와 forward/reverse composition을 검사한다.
\end{enumerate}

\newpage

\subsection{실행해서 무엇을 볼
것인가}\label{uxc2e4uxd589uxd574uxc11c-uxbb34uxc5c7uxc744-uxbcfc-uxac83uxc778uxac00}

패키지 루트에서 실행한다. \texttt{-\/-plan}은 경로만 확인하며 계산
성공을 뜻하지 않는다.

\begin{verbatim}
python tools/run_assignment.py A14 --plan
python tools/run_assignment.py A14 --run \
  --workspace learner --output ../A14_result.json
\end{verbatim}

\subsection{결과를
설명하기}\label{uxacb0uxacfcuxb97c-uxc124uxba85uxd558uxae30}

이 과제의 조건은 물질적 계면을 따르는 변위와 Lagrangian 압력의 연속이다.
두 one-sided 값을 잇는 그림의 직선을 실제 층으로 읽지 않는다.

\textbf{남길 증거.} 전달식 유도; W 및 ΔP 좌우값; 세그먼트 경계와 side
기록

\textbf{판단 기준.} 2점 예: V를 무조건 연속시키면 ΔP가 어긋나는 반례와
exact-stop 기록을 함께 보인다. 최초 예측의 정오보다, 결과를 본 뒤 이유를
고치고 다른 입력에도 적용하는지를 본다.

\subsection{막혔을 때}\label{uxb9c9uxd614uxc744-uxb54c}

\textbf{힌트 1.} q\_h는 enthalpy ratio이고 q\_P=ln(Pstart/P)와 다른
변수이다.

\textbf{힌트 2.} χ-only knot에서는 state가 연속일 수 있으므로 material
jump와 구별한다.

\subsection{다음 만남에 다시
묻기}\label{uxb2e4uxc74c-uxb9ccuxb0a8uxc5d0-uxb2e4uxc2dc-uxbb3buxae30}

다음 회차에는 밀도 jump를 없앤 극한을 설명한다. 압축응답만 바뀌는 knot와
같은 전달식을 적용해야 하는지도 구별한다. 먼저 짧게 답한 뒤 본문과
자신의 기록을 열어 확인한다.

\textbf{범위.} rapid conversion, surface shell, elasticity,
zero-frequency limit는 포함하지 않는다.

\end{document}
